\section{Related Work}
\label{sec:related}

\subsection{Domain-Specific LLM Adaptation}
\label{sec:related:domain}

Fine-tuning LLMs for specialized domains has shown promising results across several fields. In medicine, Med-PaLM~\cite{singhal2023medpalm} achieved expert-level performance on medical question answering through instruction tuning on curated clinical datasets, while PMC-LLaMA~\cite{wu2024pmcllama} demonstrated that continued pre-training on biomedical literature followed by instruction tuning can substantially improve domain performance. In the legal domain, models fine-tuned on case law and statutory text have shown improved performance on legal reasoning tasks~\cite{cui2023chatlaw}. For code generation, Code Llama~\cite{roziere2024codellama} showed that domain-specific continued pre-training followed by instruction tuning outperforms general-purpose models.

A common finding across these efforts is that domain adaptation is most effective when high-quality, domain-specific training data is available. However, for niche technical fields like reliability engineering where experts are proficient in multiple technical disciplines such as probability theory, statistics, and systems engineering, no such datasets or benchmarks currently exist. Our work addresses this gap by developing both a data generation pipeline and a fine-tuning methodology specifically designed for low-resource technical domains.

\subsection{Self-Instruct and Synthetic Data Generation}
\label{sec:related:selfinst}

Self-Instruct~\cite{wang2023selfinstruct} introduced a framework for bootstrapping instruction-following data by using a language model to generate its own training examples from a small seed set of 175 human-written tasks. The method produced 52K instruction-output pairs, improving GPT-3's instruction following by $+33.1\%$ on Super-NaturalInstructions. To ensure diversity, generated instructions are only added to the task pool if their ROUGE-L similarity with all existing instructions falls below 0.7. However, this diversity filter says little about \emph{correctness}, in fact, only 54\% of generated examples were rated as fully valid, with output quality identified as the weakest component.

Several works have extended this framework. Stanford Alpaca~\cite{taori2023alpaca} scaled self-instruct to generate 52K examples using GPT-3.5 for fine-tuning LLaMA. WizardLM~\cite{xu2023wizardlm} introduced Evol-Instruct, which iteratively increases instruction complexity through depth and breadth evolution. Ranaldi and Freitas~\cite{ranaldi2024selfrefine} proposed Self-Refine Instruction-Tuning, a two-stage approach where LLMs first generate chain-of-thought demonstrations for smaller student models, which then self-refine via DPO.

For verifiable tasks such as mathematics, answer-consistency filtering has been proposed as a quality control mechanism~\cite{yu2024cotselfinst}. The idea is straightforward: after generating a question-answer pair, the same model is prompted to solve the question independently a $n$ times; if more than $n/2$ instances agree on the anwer, the example is considered reliable and retained, otherwise it is discarded. This targets the correctness gap left by ROUGE-L, which only filters for diversity. However, a subtle issue arises when the same model is used for both generation and verification. We observe empirically that models can already solve the vast majority of questions they generate. In our domain, baseline models achieve 92--97\% accuracy on self-instruct-generated questions while scoring only 68--76\% on real textbook problems. This means same-model answer consistency is trivially satisfied and does not filter for difficulty or correctness in any meaningful way. While this does not invalidate the downstream fine-tuning results of these approaches (the trained model may still improve through exposure to diverse reasoning patterns), it reveals that answer consistency is not a reliable proxy for data quality in technical domains.

This observation is related to a broader limitation identified by SearchInstruct~\cite{li2025searchinstruct}: self-instruct ``depends solely on internal knowledge,'' meaning the generator can only create questions within its own capability envelope. They propose grounding generation in retrieved documents to overcome this.

Our approach addresses the correctness problem through \emph{cross-model} answer verification: an independent model (Claude Opus 4) generates questions from the seed dataset and solves them, we then use a comparably powerful model (ChatGPT 5.4) to verify the solution by solving each question without seegin the answer or reasoning trace, and only items where both models agree on the numerical answer (within 5\% tolerance) are retained. Because the verifier is a different model with different failure modes, agreement provides a stronger signal than same-model consistency. We additionally use a stronger generator model than the fine-tuning target, and generate full chain-of-thought triplets (question, reasoning, answer) rather than instruction-output pairs.

\subsection{Data Augmentation for Multi-Step Mathematical Reasoning}
\label{sec:related:augment}

Data augmentation has emerged as a critical strategy for improving mathematical reasoning in LLMs. MetaMath~\cite{yu2024metamath} demonstrated that rephrasing mathematical questions from multiple perspectives (without changing the underlying problem) significantly improves reasoning performance, achieving state-of-the-art results on GSM8K and MATH benchmarks. The key insight is that surface-level diversity in question formulation is more valuable than increasing problem difficulty. PersonaMath~\cite{wang2024personamath} extended this idea by generating question variants from different persona perspectives.

Ding~\etal~\cite{ding2024dataaugllm} provide a comprehensive survey of LLM-based data augmentation, finding that while larger generator models produce higher-quality data, cross-model validation is essential to prevent self-reinforcing biases. Shumailov~\etal~\cite{shumailov2024curse} formalized this concern as ``model collapse,'' showing that training on model-generated data without external grounding leads to progressive quality degradation across generations.

Our paraphrase augmentation strategy is directly inspired by MetaMath: we use Claude Opus 4.6 to rephrase questions while strictly preserving all numerical values and mathematical relationships, then verify meaning preservation with GPT-5.4. This approach scales our dataset from 280 to 866 examples and yields our strongest results ($+18.6\%$, $p=0.002$), consistent with the finding that surface diversity outperforms difficulty scaling.

\subsection{Reinforcement Learning for Reasoning}
\label{sec:related:rl}

Reinforcement learning from human feedback (RLHF)~\cite{ouyang2022training} established the paradigm of aligning LLMs with human preferences through reward modeling and policy optimization. Direct Preference Optimization (DPO)~\cite{rafailov2023dpo} simplified this by eliminating the need for a separate reward model, directly optimizing the policy from preference pairs. However, both methods rely on human-generated or model-generated preference data, which can be noisy for technical domains.

For tasks with verifiable answers, such as mathematics, reinforcement learning from verifiable rewards offers a compelling alternative. DeepSeek-Math~\cite{shao2024deepseekmath} introduced Group Relative Policy Optimization (GRPO), which generates multiple candidate responses per question, scores them against ground-truth answers, and computes group-relative advantages without requiring a critic model. DeepSeek-R1~\cite{deepseek2025r1} scaled this approach, demonstrating that GRPO can elicit sophisticated reasoning behaviors including self-verification and chain-of-thought refinement.

Reliability engineering problems are particularly well-suited for GRPO because they produce verifiable numerical answers that can serve as ground-truth rewards. In our pipeline, SFT first provides domain knowledge and reasoning format, after which GRPO optimizes the model to prefer reasoning chains that lead to correct final answers. Our early experiments with GRPO on small datasets (215 questions) showed limited improvement over SFT alone ($+0.5\%$ vs.\ $+2.3\%$), motivating the use of larger augmented datasets for the RL stage.

\TODO{Add more on GRPO with our results form ELORA.}

\subsection{Parameter-Efficient Fine-Tuning}
\label{sec:related:peft}

Low-Rank Adaptation (LoRA)~\cite{hu2022lora} enables efficient fine-tuning by adding trainable low-rank matrices to frozen pre-trained weights, drastically reducing the number of trainable parameters. QLoRA~\cite{dettmers2023qlora} combines LoRA with 4-bit quantization, enabling fine-tuning of large models on consumer-grade hardware. NEFTune~\cite{jain2023neftune} demonstrated that adding uniform noise to embedding vectors during fine-tuning acts as an effective regularizer, improving instruction-following performance by 25--35\% on standard benchmarks.

We employ 4-bit QLoRA with LoRA applied to all attention and MLP projection layers, combined with NEFTune ($\alpha=5$) as a regularizer against overfitting on our small training sets (approximately 780 examples per fold in 10-fold cross-validation on 866 total examples).
